% Based on Section 4.2 of Completing the Chain (TCHES 2026).
\begin{frame}{Functional Correctness: Proof Approach}
  \begin{block}{Goal}
    Prove that the Jasmin implementation of \xmssmt agrees with the
    RFC-like EasyCrypt specification from Section~3 of the paper.
  \end{block}
  \vspace{2mm}
  \begin{itemize}
    \setlength{\itemsep}{3mm}
    \item Extract an EasyCrypt representation of the Jasmin program
      using the Jasmin compiler.
    \item Prove \textbf{pRHL equivalence lemmas} relating the extracted
      code to the specification, under matching parameters and inputs.
    \item Show that both sides make consistent calls to the underlying
      hash function; assume correctness of the SHA-2 implementation.
  \end{itemize}
\end{frame}

\begin{frame}{Functional Correctness: Relating Representations}
  \begin{itemize}
    \setlength{\itemsep}{4mm}
    \item \textbf{Specification:} records for keys and signatures.
    \item \textbf{Implementation:} byte arrays and memory buffers.
    \item Encoding and decoding operators connect these representations
  \end{itemize}
  \vspace{3mm}
  \begin{block}{What the equivalence lemmas establish}
    \textbf{Key generation:} implementation key buffers correspond to
    the specification's public and secret keys.\par\vspace{2mm}
    \textbf{Verification:} the implementation returns a 64-bit zero
    \emph{if and only if} the specification accepts the signature.
  \end{block}
\end{frame}

\begin{frame}{Functional Correctness: Reasoning About Memory}
  \begin{itemize}
    \setlength{\itemsep}{3mm}
    \item Raw pointers require explicit reasoning about address arithmetic
      and the contents of memory.
    \item \textbf{\texttt{valid\_ptr\_i}:} adding an accessed non-negative
      offset does not overflow the 64-bit address space.
    \item \textbf{\texttt{disjoint\_ptr}:} the specified memory regions
      do not overlap.
  \end{itemize}
  \vspace{3mm}
  \begin{block}{Example: signature verification}
    The signed-message input, recovered-message output, and output-length
    buffer are assumed disjoint. Writes to the outputs cannot modify
    input data or interfere with each other.
  \end{block}
\end{frame}

\begin{frame}{Functional Correctness: TreeHash}
  \begin{itemize}
    \setlength{\itemsep}{3mm}
    \item \textbf{Most demanding proof:} relate the intermediate stacks
      and loop state of the implementation and specification.
    \item The relational loop invariant establishes \textbf{identical
      stack accesses} and matching results, avoiding a second proof
      of stack bounds at this stage.
    \item \textbf{Safety checked separately:} Jasmin's safety checker
      guarantees that the implementation performs no out-of-bounds access.
    \item Stack accesses in the RFC specification are proved within bounds
      later, when connecting it to the security proof.
  \end{itemize}
  \vspace{2mm}
  \begin{block}{Proof structure}
    The equivalence proof establishes matching computation;
    the separate safety check establishes safe memory accesses.
  \end{block}
\end{frame}
